\section{Introduction}
\subsection{Problem and motivation}
Commercial following carts reduce the need to push golf equipment but do not inherently provide bag inventory or green-distance information. Separate golf-tracking systems provide some of these functions, often through additional wearable or mobile devices \cite{stewart,arccos,arccoswatch}. Prior work surveys person-following technologies and their tradeoffs \cite{hall2025golfGps}; field studies highlight measurement limitations in golf GPS tracking \cite{stewartQFollow}. ASCC combines transport and on-cart information rather than attempting automatic shot logging.
\subsection{Objective, scope and operational scenario}
The stakeholder is a walking golfer who wants a phone-free round. The system boundary includes the motorized chassis, control electronics, sensors, tagged-club monitoring, GPS and display. External interactions include a golfer carrying a tracking device, terrain and obstacles, satellite signals, green-coordinate data, and the operator's wireless emergency stop. The planned demonstration uses one golfer, one bag, open dry grass and supervised operation. The cart reports distance from its own position, not the ball or golfer. Wet surfaces, steep slopes, crowded areas and autonomous course navigation are excluded.
\subsection{System context}
\begin{figure}[ht]\centering
\begin{tikzpicture}[>=Latex,node distance=15mm,every node/.style={font=\small},box/.style={draw,rounded corners,align=center,minimum width=27mm,minimum height=10mm}]
\node[box,fill=gray!15] (core) {ASCC system};
\node[box,left=of core] (golfer) {Golfer / beacon};
\node[box,above=of core] (course) {GPS / green data};
\node[box,right=of core] (bag) {Tagged clubs};
\node[box,below=of core] (terrain) {Terrain / obstacles};
\draw[<->] (golfer)--node[above]{range / motion}(core);
\draw[->] (course)--(core);\draw[->] (bag)--(core);\draw[->] (terrain)--(core);
\end{tikzpicture}\caption{ASCC context and external information flows (original schematic).}\end{figure}
\subsection{Scope exclusions}
Shot classification, club recommendations, mobile-app integration and autonomous navigation to waypoints are stretch possibilities only; none are required for MVP acceptance.\par
